% !TEX TS-program = XeLaTeX

%%%%%%%%%%%%%%%%%%%%%%%%%%%%%%%%%%%%%%%%%
% Dictionary
% LaTeX Template
% Version 1.0 (20/12/14)
%
% This template has been downloaded from:
% http://www.LaTeXTemplates.com
%
% Original author:
% Vel (vel@latextemplates.com) inspired by a template by Marc Lavaud
%
% License:
% CC BY-NC-SA 3.0 (http://creativecommons.org/licenses/by-nc-sa/3.0/)
%
%%%%%%%%%%%%%%%%%%%%%%%%%%%%%%%%%%%%%%%%%

% Modified by Fryderyk Mazurek

%----------------------------------------------------------------------------------------
%	PACKAGES AND OTHER DOCUMENT CONFIGURATIONS
%----------------------------------------------------------------------------------------

\documentclass[10pt,a4paper,twoside,titlepage]{report} % 10pt font size, A4 paper and two-sided margins

\usepackage{polski}
\usepackage[top=2.0cm,bottom=1.5cm,left=1.5cm,right=1.5cm,columnsep=15pt]{geometry} % Document margins and spacings
\usepackage{palatino} % Use the Palatino font
\usepackage{microtype} % Improves spacing
\usepackage{multicol} % Required for splitting text into multiple columns
\usepackage[bf,sf,center]{titlesec} % Required for modifying section titles - bold, sans-serif, centered
\usepackage{fancyhdr} % Required for modifying headers and footers
\usepackage{xeCJK}
\usepackage[hidelinks, linktoc=all]{hyperref}
\usepackage{setspace}
\usepackage{enumitem}
\usepackage{tocloft}

\setCJKmainfont{BabelStoneHan}

\fancyhead[L]{\textsf{\rightmark}} % Top left header
\fancyhead[R]{\textsf{\leftmark}} % Top right header

\renewcommand{\headrulewidth}{1.4pt} % Rule under the header

\fancyfoot[C]{\textbf{\textsf{\thepage}}} % Bottom center footer

\setlength{\headheight}{12.2pt}

\renewcommand{\footrulewidth}{1.4pt} % Rule under the footer

\newcommand*\circled[1]{\textcircled{\tiny{#1}}}

\pagestyle{fancy} % Use the custom headers and footers throughout the document

\setCJKmainfont[AutoFakeBold=3.0]{BabelStoneHan}
\newcommand*\CJKfakebold[1]{#1} % ta definicja jest potrzebna tylko po to, aby nie zmieniac wpisow, normalnie pogrubienie robi sie teraz samo

\setlength{\headheight}{12.3pt}
\addtolength{\topmargin}{-0.1pt}
\setlength{\cftsecnumwidth}{3.0em}

\begin{document}

\input{dictionary_title.tex}

\begin{titlepage}
	\centering
	\vspace*{6.5cm}
	\dictionaryTitle
	\vspace{2cm}
	{\Large\itshape Fryderyk Mazurek\par}
	
	\vfill

% Bottom of the page
	{\large Stan na dzień \today\par}
\end{titlepage}

\tableofcontents

\input{dictionary_japanese_index.tex}
\input{dictionary_polish_index.tex}
%\input{dictionary_entries.tex}
\input{dictionary_jmdict_entries.tex}


\newpage
\section*{Informacje dodatkowe}
\markboth{}{}

\noindent Słownik udostępniona na licencji \textit{Creative Commons Attribution-NonCommercial-ShareAlike 4.0 International License}. Treść licencji znajduje się na stronie \url{https://creativecommons.org/licenses/by-nc-sa/4.0/}.\\ \\
\noindent Słownik wykorzystuje dane udostępnione przez Electronic Dictionary Research And Development Group na zasadach opisanych na stronie \url{https://www.edrdg.org/edrdg/licence.html}. \\ \\
\noindent Cześć informacji zawartych w słowniku pochodzi z Wikipedii udostępnionej na licencji \textit{Creative Commons Attribution-ShareAlike 4.0 International}. Treść licencji znajduje się na stronie \url{https://creativecommons.org/licenses/by-sa/4.0/}. \\ \\
\noindent \url{https://www.japonski-pomocnik.pl} \\ \\
\noindent \url{https://play.google.com/store/apps/details?id=pl.idedyk.android.japaneselearnhelper}

%------------------------------------------------
\end{document}
